\subsection{Transformer and Ensemble Results}

Transformer models are evaluated across RoBERTa and DeBERTa architectures
with input lengths of 256 and 512 tokens. On the development split, no-seed
models provide a consistent basis for comparison across configurations.

Across all settings, performance converges within three fine-tuning epochs,
with peak macro-F1 typically reached at the second or third epoch, indicating
stable and fast convergence. Multiple random initializations are trained to
reduce variance and support robust ensembling.

Final predictions are obtained via an ensemble selected based on
development-set performance and stability, and this ensemble is used for
submission.

\begin{table}[H]
\centering
\caption{Best no-seed transformer models on the development split. For each
configuration, the highest macro-F1 score and the corresponding training
epoch are reported.}
\label{tab:transformer_best_dev}
\begin{tabular}{lccc}
\hline
\textbf{Model} & \textbf{Max Len} & \textbf{Best Epoch} & \textbf{Macro-F1 (DEV)} \\
\hline
DeBERTa & 256 & 2 & 0.743 \\
DeBERTa & 512 & 2 & 0.747 \\
RoBERTa & 256 & 3 & 0.749 \\
RoBERTa & 512 & 3 & \textbf{0.756} \\
\hline
\end{tabular}
\end{table}
